
This paper investigated whether the effective rank of pretrained weight matrices predicts per-layer optimal LoRA rank without any task data or training. For BERT-base-uncased, erank($W_0$) correlates with PARA oracle ranks at Pearson $r = 0.984$ ($p = 0.0013, n = 5$ oracle-measured layers). FFN intermediate matrices (erank $\approx$ 704--710) receive oracle rank $r^* = 64$; attention layers --- including output projections and a query projection (erank $\approx$ 557--590) --- receive oracle rank $r^* = 4$. The erank metric separates these two structural tiers without observing any training example. The participation ratio agrees with erank at $\rho = 0.968$.

Three limitations bound the scope of these findings: (1) the sample covers 5 of 72 BERT target layers, all exhibiting bimodal oracle structure; (2) only one of three planned model families has been evaluated; and (3) the coefficient of variation of erank ($\approx 0.04)$ falls below the pre-registered A1 threshold for NLP encoders. The pre-registered multi-family gate ($\geq 2/3$ families satisfying $r \geq 0.65$) has not been met.

The primary contribution is the first empirical test of a purely structural per-layer LoRA rank predictor, together with the PARA oracle + erank correlation protocol for replication and extension. Three next steps follow directly from the current limitations: (1) full oracle sweeps for BERT-base-uncased (covering intermediate erank values), DeBERTa-v3-base, and ViT-base-patch16-224; (2) partial correlation analysis controlling for layer depth to separate erank signal from depth confound; and (3) task-agnosticity assessment (H-M3: Spearman $\rho$ between MNLI and SST-2 oracle ranks for DeBERTa-v3-base).

